% Generated draft: requires mathematical and source review.
\documentclass[11pt]{article}
\usepackage[margin=25mm]{geometry}
\usepackage{fontspec}
\setmainfont{Latin Modern Roman}
\usepackage{amsmath,amssymb,mathtools,unicode-math}
\setmathfont{Latin Modern Math}
\usepackage{enumitem,needspace,xurl,hyperref}
\setlength{\parindent}{0pt}
\setlength{\parskip}{5pt}
\setlength{\emergencystretch}{3em}
\newcommand{\N}{\mathbb{N}}
\newcommand{\Z}{\mathbb{Z}}
\newcommand{\Q}{\mathbb{Q}}
\newcommand{\R}{\mathbb{R}}
\newcommand{\C}{\mathbb{C}}
\newcommand{\F}{\mathbb{F}}
\newcommand{\A}{\mathbb{A}}
\newcommand{\PP}{\mathbb{P}}
\newcommand{\E}{\mathbb{E}}
\newcommand{\eps}{\varepsilon}
\newcommand{\dd}{\,\mathrm{d}}
\begin{document}
\section*{UnsolvedMath ID: 20001852}
\textbf{Source problem code:} catalogue entry
\subsection{Definitions and mathematical statement}
The source, attributed to Farber/Panina, asks for the symplectic-volume proportion of unknots among closed spatial polygons with prescribed edge lengths. Its extracted numerator is damaged and its equality-case remark is degenerate. The following reconstruction fixes the conventional moduli problem explicitly.

For \(n\geq3\) and \(\ell=(\ell_1,\ldots,\ell_n)\in(0,\infty)^n\), assume strict feasibility \(\max_i\ell_i<\frac12\sum_i\ell_i\) and genericity
\[
 \sum_i\sigma_i\ell_i\ne0\qquad\text{for every }\sigma_i\in\{\pm1\}.
\]
The oriented, labelled polygon space modulo translations and orientation-preserving rotations is
\[
 M_\ell=\{(e_1,\ldots,e_n)\in(\R^3)^n:
 |e_i|=\ell_i,\ \sum_i e_i=0\}/SO(3).
\]
It is the symplectic reduction of the product of coadjoint two-spheres, with symplectic form \(\omega\) and dimension \(2n-6\). Put \(r=n-3\) and use Liouville measure \(\omega^r/r!\). Let \(U_0\subset M_\ell\) be the locus where the associated polygonal closed curve is embedded and ambient isotopic to a standard circle in \(\R^3\). The intended proportion is
\[
 P_0(\ell)=\frac{\operatorname{Vol}_\omega(U_0)}
 {\operatorname{Vol}_\omega(M_\ell)}.
\]
If the extracted \(U_i\) denote connected knot-type strata, the numerator means their union over unknot strata. Self-intersecting polygons are excluded from \(U_0\); in the generic regular setting their discriminant has zero Liouville measure. The original text does not define its strata or volume convention.

The supplied aid states that every feasible generic vector with \(3\leq n\leq5\) has \(P_0(\ell)=1\). It further states
\[
 \exists\epsilon_0>0\ \forall\epsilon\in\R:\quad
 0<|\epsilon|<\epsilon_0\Longrightarrow
 \ell_\epsilon=(1,1,1,1,1,1+\epsilon)
\]
is positive, feasible and generic, and both its embedded unknot locus and its embedded trefoil locus contain open subsets of positive symplectic volume. Consequently
\[
 0<P_0(\ell_\epsilon)<1.
\]
These are reported partial results, not a formula for arbitrary length vectors.

A precise action--angle form of the aid's reduction uses fan diagonals. With vertices \(v_0,\ldots,v_n=v_0\), set \(d_j=|v_j-v_0|\), \(d_1=\ell_1\), \(d_{n-1}=\ell_n\), and let \(\Delta_\ell\subset\R^r\) be the admissible region for \(d_2,\ldots,d_{n-2}\) defined by
\[
 |d_{j-1}-\ell_j|\leq d_j\leq d_{j-1}+\ell_j
 \quad(2\leq j\leq n-1).
\]
On the full-measure nondegenerate action--angle locus, let \(\Phi(d,\theta)\) reconstruct the polygon, with each bending angle of period \(2\pi\). In the convention \(\omega=\sum\dd d_j\wedge\dd\theta_j\),
\[
 P_0(\ell)=\frac{
 \displaystyle\int_{\Delta_\ell}\int_{(\R/2\pi\Z)^r}
 \mathbf1_{U_0}(\Phi(d,\theta))\,\dd\theta\,\dd d}
 {(2\pi)^r\operatorname{Vol}_r(\Delta_\ell)}.
\]
A common normalization factor in the symplectic form cancels in the ratio. This is an exact integral reduction, not an evaluation of its knot-type indicator.

Finally, the source's condition \(\ell_1+\cdots+\ell_{n-1}=\ell_n\) forces equality in the vector triangle inequality: the first \(n-1\) edges align and the last returns along the same line. This is a degenerate collinear configuration, not an embedded unknot. It lies outside the generic positive-volume setup above, so it does not justify the source's asserted probability-one equality case.
\subsection{Short English statement}
The problem asks how much of a prescribed-length polygon space consists of unknots, measured by its symplectic volume. The aid reports proportion one for at most five edges and a proportion strictly between zero and one for small generic perturbations of equilateral hexagons. The source's equality-of-lengths example is collinear and degenerate.
\subsection{Sources}
\begin{itemize}
\item Permanent UnsolvedMath ID \(20001852\), exact public code catalogue entry; local registry \texttt{aim-geometry-notes.json}, \texttt{notes[189]}, canonical tag problem.
\item AIM workshop ``Configuration spaces of linkages,'' Geometry, source item 3, attributed in the extracted text to Farber/Panina: \url{https://aimath.org/pastworkshops/linkagesproblems.pdf}. The URL is retained as supplied; no PDF was accessed for this draft.
\item The source attributes the symplectic structure to Klyatchko, without a full citation. The original literature field is empty and no proposer/year metadata are supplied beyond the in-text attribution.
\item Research attempt 1, ``The symplectic-volume proportion of unknotted spatial polygons,'' labelled partial\_result, partially\_solved, proved and low novelty confidence. Although the statement-status field says exact, its verification note explicitly describes the extracted text as visibly damaged and identifies zero-based registry index 189.
\end{itemize}
\end{document}
